\begin{table*}[t]
\centering\small
\setlength{\tabcolsep}{4pt}
\renewcommand{\arraystretch}{1.09}
\begin{tabularx}{\textwidth}{@{}>{\raggedright\arraybackslash}p{.17\textwidth}>{\raggedright\arraybackslash}X>{\raggedright\arraybackslash}X@{}}
\toprule
Object & Study A: paired V5 & Study B: complete native corpus \\
\midrule
Knowledge & Source-extracted propositions, macro graph, refined local evidence rules and mapped routes & Common deterministic compilation of all three public tenancy templates \\
Online gate & Rule-local true/false/unresolved guard; macro edges do not gate requests & Inherited applicability scopes, obligation condition and separate acquisition permission \\
Evidence state & Held document types; true-guard confirmation; union of mapped routes & Presence, capability/route adequacy, multi-document joint adequacy; per-capability route selection \\
Population & 9 development + 27 held-out pairs; nine shared concepts, four contexts; 72 units total & 60 development/11 families + 90 protected/17 families; 150 claims, three tenancy domains \\
Endpoint & Route-aware signed additions/withdrawals, 27-pair micro scores & Complete native artifact and separate actual immediate-checklist metrics \\
\bottomrule
\end{tabularx}
\caption{A common source-to-state-to-evidence interface does not make the two
implementations, populations or measurements interchangeable. Study B extends
the evidence question rather than retrospectively re-evaluating V5.}
\label{tab:correspondence}
\end{table*}
